% !TEX program = pdflatex
% Paper 2 (v2): exact redundancy of degree-based topological indices and a
% target-aware screen for new molecular descriptors.
% Compile: pdflatex (twice).  Figures are read from ../figures/ or the
% bundle directory.

\documentclass[11pt,a4paper]{article}

\usepackage[T1]{fontenc}
\usepackage[utf8]{inputenc}
\usepackage{lmodern}
\usepackage[margin=1in]{geometry}
\usepackage{amsmath,amssymb,amsthm}
\usepackage{mathtools}
\usepackage{booktabs}
\usepackage{array}
\usepackage{graphicx}
\graphicspath{{../figures/}{./}}
\PassOptionsToPackage{hyphens}{url}
\usepackage{hyperref}
\hypersetup{colorlinks=true,linkcolor=blue!50!black,citecolor=blue!50!black,urlcolor=blue!50!black,breaklinks=true}
\IfFileExists{xurl.sty}{\usepackage{xurl}}{}
\usepackage{xcolor}
\usepackage{enumitem}
\setlist[itemize]{topsep=2pt,itemsep=2pt}
\setlist[enumerate]{topsep=2pt,itemsep=2pt}
\usepackage{placeins}
\usepackage{float}
\newcommand{\oracleRmseEsol}{0.85}
\newcommand{\oracleRsqEsol}{0.83}
\newcommand{\looRmseEsol}{1.87}
\newcommand{\oracleRmseFsv}{2.53}
\newcommand{\oracleRsqFsv}{0.57}
\newcommand{\looRmseFsv}{3.87}
\newcommand{\oracleRmseLipo}{0.26}
\newcommand{\oracleRsqLipo}{0.95}
\newcommand{\looRmseLipo}{1.11}
\newcommand{\oracleAucBbbp}{0.999}
\newcommand{\looAucBbbp}{0.62}
\newcommand{\rdkitRsqEsol}{0.90}
\newcommand{\rdkitRsqFsv}{0.89}
\newcommand{\rdkitRsqLipo}{0.65}
\newcommand{\nGenDatasets}{13}
\newcommand{\nGenCertified}{12}
\newcommand{\nGenMolecules}{318\,648}
\newcommand{\hivHyper}{350}
\newcommand{\hivPairs}{37}
\newcommand{\hivRankM}{36}
\newcommand{\hivRankB}{17}
\newcommand{\hivRestrictedN}{40\,777}
\newcommand{\hivRestrictedResid}{5.0\times10^{-12}}
\newcommand{\genMaxResid}{1.6\times10^{-9}}
\newcommand{\nullBandB}{0.133}
\newcommand{\nullBandA}{0.102}
\newcommand{\nullKB}{60}
\newcommand{\nullKA}{100}
\newcommand{\nullPmin}{0.016}
\newcommand{\nullBEsol}{0.12}
\newcommand{\nullAEsol}{0.08}
\newcommand{\nullBFsv}{0.09}
\newcommand{\nullAFsv}{0.10}
\newcommand{\nullBLipo}{0.08}
\newcommand{\nullALipo}{0.05}
\newcommand{\nullBBbbp}{0.13}
\newcommand{\nullABbbp}{0.05}
\newcommand{\nullSurvivors}{0}
\newcommand{\qInfoEsol}{0.437}
\newcommand{\qFourBbbp}{0.437}
\newcommand{\chemPassTotal}{0}
\newcommand{\chemTopoBeyond}{1}
\newcommand{\chemCvGainMax}{0.024}

\usepackage{xr}
\externaldocument{paper2_orthogonality_screening}
\renewcommand{\thesection}{S\arabic{section}}
\renewcommand{\thetable}{S\arabic{table}}
\renewcommand{\thefigure}{S\arabic{figure}}
\setlength{\emergencystretch}{2em}

\newtheorem{theorem}{Theorem}[section]
\newtheorem{proposition}[theorem]{Proposition}
\newtheorem{lemma}[theorem]{Lemma}
\newtheorem{corollary}[theorem]{Corollary}
\theoremstyle{definition}
\newtheorem{definition}[theorem]{Definition}
\newtheorem{remark}[theorem]{Remark}

\newcommand{\pcor}{\mathrm{pcor}}
\newcommand{\Gc}{\mathcal{G}}
\newcommand{\col}{\operatorname{col}}

\title{Supplementary Information for:\\ Ten Edge-Degree Counts Are All a Degree-Based Topological Index Can See}

\author{
  Ganesh Shiwakoti\thanks{Department of Mathematics and Statistics, Florida Atlantic University, USA. \texttt{ganshiwakoti@gmail.com}} \and
  Advik Natarajan\thanks{Department of Mathematics, Loyola College, Chennai, India.} \and
  Michael Arockiaraj\footnotemark[2]
}

\date{Draft v2.1 of \today}

\begin{document}
\maketitle

\section{The baseline}\label{sup:baseline}
\begin{table}[H]
\centering
\small
\caption{The $30$-index classical baseline.}
\label{tab:baseline}
\begin{tabular}{lp{0.68\linewidth}}
\toprule
Family & Indices\\
\midrule
Degree-based, BID (18) & $M_1$, $M_2$, $mM_1$, $mM_2$, $F$, $R$, SCI, $H$, $GA$~\cite{VukicevicFurtula2009GA}, $AG$, ABC, ABS, AZI, $SO$~\cite{Gutman2021Sombor}, $SO_{\mathrm{red}}$, Alb~\cite{Albertson1997}, $\sigma$, $\mathrm{redM}_1$\\
Distance-based (8) & $W$, WW, HR, degree distance, Gutman, Mostar, $Sz$, $PI$\\
Spectral (3) & Estrada EE, energy, spectral radius\\
Other (1) & Balaban $J$\\
\bottomrule
\end{tabular}
\end{table}

\section{The value matrix of the baseline}\label{sup:phi}
\begin{table}[h]
\centering
\caption{The value matrix $\Phi_B$ of the $18$ degree-based baseline indices on the ten degree pairs $(i,j)\in P_4$: entry $f_b(i,j)$, obtained by evaluating each index on $K_{i,j}$ and dividing by its $ij$ edges. Vertex-degree sums $\sum_v g(d_v)$ are written as $f(a,b)=g(a)/a+g(b)/b$. The matrix has rank $10$ (smallest/largest singular value $1.2e-04$).}
\label{tab:phi}
\resizebox{\textwidth}{!}{%
\begin{tabular}{llrrrrrrrrrr}
\toprule
Index & $f(a,b)$ & $(1,1)$ & $(1,2)$ & $(1,3)$ & $(1,4)$ & $(2,2)$ & $(2,3)$ & $(2,4)$ & $(3,3)$ & $(3,4)$ & $(4,4)$\\
\midrule
$M_1$ & $a+b$ & $2$ & $3$ & $4$ & $5$ & $4$ & $5$ & $6$ & $6$ & $7$ & $8$\\
$M_2$ & $ab$ & $1$ & $2$ & $3$ & $4$ & $4$ & $6$ & $8$ & $9$ & $12$ & $16$\\
$mM_1$ & $a^{-3}+b^{-3}$ & $2$ & $1.12$ & $1.04$ & $1.02$ & $0.25$ & $0.162$ & $0.141$ & $0.0741$ & $0.0527$ & $0.0312$\\
$mM_2$ & $1/(ab)$ & $1$ & $0.5$ & $0.333$ & $0.25$ & $0.25$ & $0.167$ & $0.125$ & $0.111$ & $0.0833$ & $0.0625$\\
F & $a^2+b^2$ & $2$ & $5$ & $10$ & $17$ & $8$ & $13$ & $20$ & $18$ & $25$ & $32$\\
R & $(ab)^{-1/2}$ & $1$ & $0.707$ & $0.577$ & $0.5$ & $0.5$ & $0.408$ & $0.354$ & $0.333$ & $0.289$ & $0.25$\\
SCI & $(a+b)^{-1/2}$ & $0.707$ & $0.577$ & $0.5$ & $0.447$ & $0.5$ & $0.447$ & $0.408$ & $0.408$ & $0.378$ & $0.354$\\
H & $2/(a+b)$ & $1$ & $0.667$ & $0.5$ & $0.4$ & $0.5$ & $0.4$ & $0.333$ & $0.333$ & $0.286$ & $0.25$\\
GA & $2\sqrt{ab}/(a+b)$ & $1$ & $0.943$ & $0.866$ & $0.8$ & $1$ & $0.98$ & $0.943$ & $1$ & $0.99$ & $1$\\
AG & $(a+b)/(2\sqrt{ab})$ & $1$ & $1.06$ & $1.15$ & $1.25$ & $1$ & $1.02$ & $1.06$ & $1$ & $1.01$ & $1$\\
ABC & $\sqrt{(a+b-2)/(ab)}$ & $0$ & $0.707$ & $0.816$ & $0.866$ & $0.707$ & $0.707$ & $0.707$ & $0.667$ & $0.645$ & $0.612$\\
ABS & $\sqrt{(a+b-2)/(a+b)}$ & $0$ & $0.577$ & $0.707$ & $0.775$ & $0.707$ & $0.775$ & $0.816$ & $0.816$ & $0.845$ & $0.866$\\
AZI & $\bigl(ab/(a+b-2)\bigr)^3\;(0\text{ on }K_2)$ & $0$ & $8$ & $3.38$ & $2.37$ & $8$ & $8$ & $8$ & $11.4$ & $13.8$ & $19$\\
SO & $\sqrt{a^2+b^2}$ & $1.41$ & $2.24$ & $3.16$ & $4.12$ & $2.83$ & $3.61$ & $4.47$ & $4.24$ & $5$ & $5.66$\\
SO$_{\mathrm{red}}$ & $\sqrt{(a-1)^2+(b-1)^2}$ & $0$ & $1$ & $2$ & $3$ & $1.41$ & $2.24$ & $3.16$ & $2.83$ & $3.61$ & $4.24$\\
Alb & $|a-b|$ & $0$ & $1$ & $2$ & $3$ & $0$ & $1$ & $2$ & $0$ & $1$ & $0$\\
$\sigma$ & $(a-b)^2$ & $0$ & $1$ & $4$ & $9$ & $0$ & $1$ & $4$ & $0$ & $1$ & $0$\\
redM$_1$ & $(a-1)^2/a+(b-1)^2/b$ & $0$ & $0.5$ & $1.33$ & $2.25$ & $1$ & $1.83$ & $2.75$ & $2.67$ & $3.58$ & $4.5$\\
\bottomrule
\end{tabular}}
\end{table}


\section{Definitions of the 20 alternative candidate indices}\label{app:candidates}
For reproducibility we give the formal definition of each of the
twenty candidate indices screened in Section~\ref{sec:nonbid}. Let
$G = (V, E)$ be a simple connected hydrogen-suppressed molecular
graph with $n = |V|$ vertices, $m = |E|$ edges, degree sequence
$\{d_v\}_{v\in V}$, geodesic distance matrix $D = (d_{uv})$,
adjacency matrix $A$, eccentricity $\varepsilon(v) = \max_u d_{uv}$,
and adjacency-spectrum / Laplacian-spectrum eigenvalues
$\{\lambda_i\}$ / $\{\mu_i\}$ respectively. All twenty definitions
are implemented in
\texttt{src/novel\_candidates.py} of the companion repository.

\paragraph{Information-theoretic (5).}
With $H(\{c_i\}) = -\sum_i (c_i / \sum_j c_j)\log(c_i / \sum_j c_j)$
the Shannon entropy of a count distribution:
\begin{align*}
\mathrm{InfoH\_deg}(G) &= H\bigl(\{|\{v : d_v = k\}|\}_k\bigr), \\
\mathrm{InfoH\_dist}(G) &= H\bigl(\{|\{(u,v) : d_{uv} = k,\, u<v\}|\}_k\bigr), \\
\mathrm{InfoH\_spec}(G) &= -\sum_{i} p_i \log p_i,\ p_i = \lambda_i^2 / \textstyle\sum_j \lambda_j^2, \\
\mathrm{InfoH\_ecc}(G) &= H\bigl(\{|\{v : \varepsilon(v) = k\}|\}_k\bigr), \\
\mathrm{Bonchev\_Id}(G) &= W\log_2 W - \sum_{k \ge 1} k\,d_k \log_2(k\,d_k),
\end{align*}
where in the last line $d_k = |\{(u,v) : d_{uv}=k,\,u<v\}|$ and
$W = \sum_k k\,d_k$ is the Wiener index; this Bonchev-type form, as
implemented, differs from the standard information-theoretic Wiener
index $I_D^W = W\log_2W-\sum_k k\,d_k\log_2 k$.

\paragraph{Centrality-based (5).}
With NetworkX centralities (betweenness normalised by $\binom{n-1}{2}$,
NetworkX's default):
\begin{align*}
\mathrm{Btw\_sum}(G) &= \sum_{v \in V} c_B(v),\quad c_B(v) = \binom{n-1}{2}^{-1}\!\!\sum_{s \ne t \ne v}\!\!
\frac{\sigma_{st}(v)}{\sigma_{st}}, \\
\mathrm{Cls\_sum}(G) &= \sum_{v} (n-1)\Bigl(\sum_{u \ne v} d_{uv}\Bigr)^{-1}, \\
\mathrm{Eig\_sum}(G) &= \sum_{v} x_v,\ x \text{ the (positive) Perron eigenvector of } A, \\
\mathrm{Harm\_cent}(G) &= \sum_{v}\sum_{u \ne v}\frac{1}{d_{uv}} = 2\,\mathrm{HR}(G), \\
\mathrm{Btw\_var}(G) &= \mathrm{Var}_v\bigl(c_B(v)\bigr).
\end{align*}

\paragraph{Spectral-hybrid (3).}
With $\mathrm{diam}(G) = \max_{u,v} d_{uv}$ and Estrada index
$\mathrm{EE}(G) = \sum_i e^{\lambda_i}$:
\begin{align*}
\mathrm{AlgConn}(G) &= \mu_2(G), \text{ second-smallest Laplacian eigenvalue}, \\
\mathrm{Rho}\!\times\!\mathrm{Dia}(G) &= \rho(G)\cdot\mathrm{diam}(G),\ \rho(G) = \max_i|\lambda_i|, \\
\mathrm{EE}\!\times\!\mathrm{W}(G) &= \log\bigl(1+\mathrm{EE}(G)\bigr)\cdot\log\bigl(1+W(G)\bigr).
\end{align*}

\paragraph{Motif / structural (4).}
\begin{align*}
\mathrm{Triangles}(G) &= \tfrac{1}{6}\operatorname{tr}(A^3), \\
\mathrm{MeanClust}(G) &= \frac{1}{n}\sum_{v} \frac{2 t(v)}{d_v(d_v-1)},\ \,t(v) = \text{triangles at } v, \\
\mathrm{FourCyc}(G) &= \tfrac{1}{8}\bigl(\operatorname{tr}(A^4) - 2\!\sum_v d_v(d_v - 1) - 2m\bigr), \\
\mathrm{CutVerts}(G) &= |\{v \in V : G - v \text{ is disconnected}\}|.
\end{align*}

\paragraph{Eccentricity-based (3).}
\begin{align*}
\mathrm{TotEcc}(G) &= \sum_{v \in V} \varepsilon(v), \\
\mathrm{Radius}(G) &= \min_{v \in V} \varepsilon(v), \\
\mathrm{MeanEcc}(G) &= \frac{1}{n}\sum_{v \in V} \varepsilon(v).
\end{align*}

None of these twenty candidates is claimed as a newly invented
index. The information-theoretic family follows the information-theoretic descriptor tradition~\cite{TodeschiniConsonniHandbook};
the centrality family follows the standard graph-theory
literature; the motif and eccentricity families are textbook
graph invariants; the two spectral-hybrid product indices
($\mathrm{Rho}\times\mathrm{Dia}$ and $\mathrm{EE}\times\mathrm{W}$) are
constructed here as controls, to test whether product indices escape
the redundancy ceiling; they are not proposed as descriptors. They
appear here purely as screening targets to test the generality
of the target-aware screen, and the
verdict reported in Section~\ref{sec:nonbid} is a statement about
this particular candidate set under the tested baselines and
thresholds, not a claim that any individual candidate is or is
not a useful QSPR descriptor.



\section{Statistical details of the benchmark}\label{sup:stats}
The benchmark of the main text uses $r=5$ repeats of $k=5$-fold
cross-validation with folds grouped by isomorphism class of the
hydrogen-suppressed graph. For a paired difference $d_{ij}$ (repeat $i$,
fold $j$) between two models, the Nadeau--Bengio corrected variance of
the mean difference is $\hat\sigma^2\,(1/(kr)+n_{\mathrm{test}}/n_{\mathrm{train}})$,
where $\hat\sigma^2$ is the sample variance of the $kr$ differences; the
corrected $t$ statistic has $kr-1$ degrees of freedom and the reported
$90\%$ intervals are $\bar d\pm t_{0.95,kr-1}\,\hat\sigma_{\mathrm{corr}}$.
Equivalence is tested by two one-sided tests (TOST) at level $0.05$
against a margin of $5\%$ of the reference model's mean RMSE, or
$0.05\,(\mathrm{AUC}_{\mathrm{ref}}-0.5)$ for BBBP; equivalence is
declared when the $90\%$ interval lies inside $[-\delta,\delta]$.
Within each dataset the five comparisons are Holm-adjusted. The
inner cross-validation that selects the LASSO or $\ell_1$-logistic
penalty is also grouped by graph. Random forests use $300$ trees.

\paragraph{Screened subsets and the combined screen.} Restricting the baseline to its pairwise-screened
subset (about $6$--$8$ indices) is equivalent within $5\%$ for the
random forest on the three regression sets and slightly better on BBBP;
for LASSO it is equivalent within the margin on the regression sets,
although the cost is distinguishable from zero on ESOL ($+0.045$) and
Lipophilicity ($+0.009$), and it is not shown equivalent on BBBP.
The combined screen, which also requires $|\pcor|\ge0.10$, is a
diagnostic and not a feature selector: on Lipophilicity it removes every
feature in $24$ of $25$ folds.

\section{Datasets used for the generality check}\label{sup:data}
\begin{table}[h]
\centering
\footnotesize
\caption{The span certificate on thirteen MoleculeNet SMILES datasets (no targets used). $n$: parsed molecules; $n_{>4}$: molecules containing an atom of degree above four; $\Delta$: largest degree; $|P_D|$: realised degree pairs; ranks of the count matrix $M$ and of the $18$-index baseline block $B$; largest relative residual of an $m_{ij}$ column on $B$; whether $\operatorname{rank}[\mathbf 1,B]=\operatorname{rank}[\mathbf 1,M]$ with residuals below $10^{-8}$.}
\label{tab:general}
\setlength{\tabcolsep}{3pt}
\resizebox{\textwidth}{!}{%
\begin{tabular}{lrrrrrrcl}
\toprule
Dataset & $n$ & $n_{>4}$ & $\Delta$ & $|P_D|$ & rk $M$ & rk $B$ & max resid. & certified\\
\midrule
ESOL & $1127$ & $0$ & $4$ & $10$ & $10$ & $10$ & $4.6\times10^{-13}$ & yes\\
FreeSolv & $639$ & $0$ & $4$ & $10$ & $10$ & $10$ & $2.1\times10^{-12}$ & yes\\
Lipophilicity & $4200$ & $0$ & $4$ & $9$ & $9$ & $9$ & $2.6\times10^{-12}$ & yes\\
BBBP & $2039$ & $0$ & $4$ & $10$ & $10$ & $10$ & $1.1\times10^{-12}$ & yes\\
BACE & $1513$ & $0$ & $4$ & $9$ & $9$ & $9$ & $1.8\times10^{-12}$ & yes\\
ClinTox & $1477$ & $1$ & $6$ & $11$ & $11$ & $11$ & $2.4\times10^{-12}$ & yes (on the $\Delta\le4$ subset, $n=1476$: ranks $10$/$10$, yes)\\
SIDER & $1396$ & $1$ & $6$ & $11$ & $11$ & $11$ & $7.7\times10^{-12}$ & yes (on the $\Delta\le4$ subset, $n=1395$: ranks $10$/$10$, yes)\\
Tox21 & $7812$ & $2$ & $6$ & $12$ & $12$ & $12$ & $6.6\times10^{-11}$ & yes (on the $\Delta\le4$ subset, $n=7810$: ranks $10$/$10$, yes)\\
ToxCast & $8566$ & $3$ & $6$ & $13$ & $13$ & $13$ & $1.6\times10^{-9}$ & yes (on the $\Delta\le4$ subset, $n=8563$: ranks $10$/$10$, yes)\\
QM8 & $21783$ & $0$ & $4$ & $10$ & $10$ & $10$ & $3.6\times10^{-11}$ & yes\\
QM9 & $133882$ & $0$ & $4$ & $10$ & $10$ & $10$ & $1.4\times10^{-10}$ & yes\\
MUV & $93087$ & $0$ & $4$ & $9$ & $9$ & $9$ & $2.0\times10^{-11}$ & yes\\
HIV & $41127$ & $350$ & $10$ & $37$ & $36$ & $17$ & $9.9\times10^{-1}$ & \textbf{no} (on the $\Delta\le4$ subset, $n=40777$: ranks $10$/$10$, yes)\\
\bottomrule
\end{tabular}}
\end{table}


The nine additional MoleculeNet files (BACE, ClinTox, SIDER, Tox21,
ToxCast, HIV, MUV, QM8, QM9) are downloaded from the DeepChem mirror by
script \texttt{78} of the repository; their SHA-256 digests are recorded
in the results file \texttt{certificate\_generality.csv}. Only the SMILES
column of each file is used.


\end{document}
